\chapter{Introducción}

\section{Propósito de esta guía}
Esta guía tiene como objetivo proporcionar las instrucciones necesarias para instalar, configurar y utilizar el \textbf{Sistema de Gestión para la Curia de Cuenca}. El sistema está diseñado para satisfacer las necesidades administrativas y de registro de la Curia, centralizando la gestión de las actividades relacionadas con la Iglesia y sus servicios, tales como bautizos, matrimonios, catequesis, pagos, ministros, parroquias e iglesias, entre otros módulos.

\section{A quiénes va dirigida}
Esta guía está pensada para dos perfiles de usuarios:
\begin{itemize}
    \item \textbf{Personal administrativo y religioso de la Curia:} utilizarán principalmente los capítulos de uso de la aplicación (Capítulo \ref{cap:uso}).
    \item \textbf{Personal técnico o encargado de instalación:} seguirán principalmente los capítulos de requisitos e instalación (Capítulos \ref{cap:requisitos} y \ref{cap:instalacion}).
\end{itemize}

\section{Alcance}
La guía cubre de forma detallada dos escenarios de despliegue:
\begin{itemize}
    \item \textbf{Escenario 1 (Servidor local):} instalación física directa sobre Windows con .NET 8 y PostgreSQL 18 para operación en red de área local (LAN).
    \item \textbf{Escenario 2 (Servidor en la Nube):} despliegue contenerizado con Docker Compose en Google Cloud Compute Engine (Ubuntu Linux) para acceso global en Internet con persistencia y respaldos automáticos.
    \item Los requisitos de hardware, software y recursos en la nube para ambos esquemas.
    \item El uso paso a paso de cada módulo de la aplicación mediante interfaces unificadas.
    \item La solución de los problemas más frecuentes en la instalación local, el contenedor de Docker y la operación diaria.
\end{itemize}

\section{Convenciones utilizadas}
A lo largo del documento se emplean los siguientes elementos visuales:
\begin{itemize}
    \item \texttt{Escriba los comandos en este estilo de letra} para comandos, rutas de archivos y código.
    \item Cuadros de \textbf{Nota}, \textbf{Consejo} y \textbf{Atención} para resaltar información importante.
    \item Los pasos numerados indican la secuencia exacta de acciones a realizar.
\end{itemize}

\begin{nota}
Las capturas de pantalla referenciadas a lo largo de este manual se encuentran alojadas en la carpeta \texttt{imagenes/} del proyecto documental.
\end{nota}
